\section{Relación entre la predicción de ruido y el \textit{score matching}}

\subsection{De la predicción de ruido a la función \textit{score}}

En las conclusiones del curso, el instructor anunció una conexión importante: existe una correspondencia profunda entre las redes de predicción de ruido $\epsilon_\theta(x_t, t)$ y la función \textit{score} en los modelos basados en \textit{score}.

La función \textit{score} se define como el gradiente de la densidad de probabilidad logarítmica respecto a los datos:

$$s_\theta(x_t, t) = \nabla_{x_t} \log p_t(x_t)$$

Se puede demostrar que existe la siguiente relación entre la red de predicción de ruido y la función \textit{score}:

$$\epsilon_\theta(x_t, t) \approx -\sqrt{1 - \bar{\alpha}_t}\, \nabla_{x_t} \log q_t(x_t)$$

\begin{knowledgebox}{Unificación de DDPM y los modelos basados en Score}
Esta conexión revela que dos marcos de modelos generativos que a primera vista parecen distintos —los modelos probabilísticos de difusión por eliminación de ruido (DDPM) y los modelos generativos basados en score (SGM)— son, en esencia, equivalentes:
\begin{itemize}
    \item DDPM aprende a predecir el ruido $\epsilon$
    \item SGM aprende a predecir la función \textit{score} $\nabla_x \log p(x)$
    \item La única diferencia entre ambos es un factor de escala relacionado con el paso temporal
\end{itemize}
Esta equivalencia fue descrita en una descripción continua unificada dentro del marco Score SDE por Song et al. (2021).
\end{knowledgebox}

\subsection{Comprensión intuitiva}

La predicción de ruido puede entenderse como responder a la siguiente pregunta: ``¿cuánto ruido hay en $x_t$?'' Por otro lado, la función \textit{score} responde: ``¿en qué dirección debería moverse $x_t$ para aumentar su densidad de probabilidad?'' En el contexto del proceso de difusión, estas dos preguntas son esencialmente equivalentes: la dirección en la que se elimina el ruido es exactamente la dirección en la que aumenta la densidad de probabilidad.

\subsection{Resumen de la sección}

Existe una correspondencia matemática precisa entre la predicción de ruido y la función \textit{score}, lo cual establece un puente entre los modelos DDPM y los basados en score, sentando las bases para el marco unificado posterior de Score SDE.

%% ========== Síntesis y ampliación ==========


